\section*{Chapter \ref{ch:s2:cross-currency-basis-and-fx-carry} --- Cross-Currency Basis and FX Carry}

\begin{solution}{exo:s2:cross-currency-basis-and-fx-carry:1}
$2.7/7.2 = 0.375$ with the rounded figures, 0.37 unrounded.
\end{solution}

\begin{solution}{exo:s2:cross-currency-basis-and-fx-carry:2}
$(50 - 15)/252 \times 10 = 1.39$ basis points of the amount lent.
\end{solution}

\begin{solution}{exo:s2:cross-currency-basis-and-fx-carry:3}
$(1.14 - 0.63)/1.30 = 0.39$ with the rounded figures, 40\% unrounded.
\end{solution}

\begin{solution}{exo:s2:cross-currency-basis-and-fx-carry:4}
Its interest rate was the lowest or near the lowest of the nine throughout, often at or below zero, so it was always among the three funding
currencies.
\end{solution}

\begin{solution}{exo:s2:cross-currency-basis-and-fx-carry:5}
It fell a little each month; a put struck 1.5 standard deviations below zero pays only on a month that falls further than that, so the premiums were
spent and little came back.
\end{solution}

\begin{solution}{exo:s2:cross-currency-basis-and-fx-carry:6}
Closing it needs balance sheet: an arbitrageur must borrow dollars directly and lend them through the swap, a gross position that banks charge for,
most of all on reporting dates. The basis settles where that charge makes the trade not worth doing for the marginal bank.
\end{solution}

\begin{solution}{exo:s2:cross-currency-basis-and-fx-carry:7}
The hedge now costs 0.55\% a year and pays 0.63\%: it makes money, and the hedged basket's Sharpe ratio rises to 0.26 with no skew. The crash premium
is not in the currencies' returns alone; it is in what option sellers charge above fair value, here the 1.2 times trailing volatility of the chapter's
pricing.
\end{solution}

\begin{solution}{exo:s2:cross-currency-basis-and-fx-carry:8}
Levering 2.8 times multiplies the crashes too: the autumn 2008 loss of 19.4\% becomes about 54\%. The Sharpe ratio hides a negative skew of $-0.60$;
size on the crash, not on the volatility.
\end{solution}

\begin{solution}{pb:s2:cross-currency-basis-and-fx-carry:1}
\begin{enumerate}
\item Long the highest-rate currencies, short the lowest; from FRED exchange rates and OECD interbank rates, monthly.
\item 2.7\% a year at 7.2\% volatility, Sharpe ratio 0.37, skewness $-0.60$; $-19.4\%$ in autumn 2008, $-10.5\%$ in October 2008.
\item Long mostly New Zealand and Australian dollars and the krone; short the Swiss franc always, the yen, the krona and the euro often.
\item Carry earned Sharpe ratios like equity factors; crash premia were at most a third of the return.
\item Carry with options that pay in a crash.
\item Unhedged 1.30\% a year, Sharpe ratio 0.23, skewness $-0.39$; hedged 0.78\%, 0.15, 0.01.
\item The first and third crashes' months were cut from $-12.0\%$ to $-7.9\%$ and from $-14.4\%$ to $-8.5\%$; the second barely changed.
\item It is what must be paid to remove the crash from the return.
\item Borrowing dollars directly or through a swap should cost the same; the basis is the gap.
\item Large, persistent deviations, strongest at quarter-ends, pointing to regulation.
\item Lend dollars through the swap; 11.2 basis points a year after a 15 basis-point cost.
\item 42.5 basis points a year while lent, over 40 days a year.
\item \textbf{Named result.} The synthetic carry basket's Sharpe ratio falls from 0.23 to 0.15 with monthly put hedges, which remove its negative skew
at 40\% of its return; the basis trade earns 42.5 basis points a year while lent in the quarter-end windows, 6.7 a year in all.
\item Forward points include the basis, and carry crashes are funding events.
\item Less carry in the calm, some protection in long crashes, none in fast ones.
\item By the crash, not the volatility.
\item Swap prices for the exact dates, and the lender's own \omterm{def:s2:swap-spread-and-asset-swap-trades:bscost}{balance-sheet cost}.
\item The basis funding arbitrage.
\item Both are prices of bank balance sheet: the swap spread for holding bonds, the basis for lending dollars.
\item The cost of bank balance sheet for lending dollars.
\end{enumerate}
\end{solution}

\begin{solution}{iq:s2:cross-currency-basis-and-fx-carry:1}
Borrow in low-rate currencies and lend in high-rate ones; exchange rates do not, on average, move enough to offset the differential, but they do in
crashes. The return is pay for bearing that crash risk.
\end{solution}

\begin{solution}{iq:s2:cross-currency-basis-and-fx-carry:2}
Hedge the carry positions with out-of-the-money options at market prices and compare hedged and unhedged returns; the difference is the premium paid
to remove the crash.
\end{solution}

\begin{solution}{iq:s2:cross-currency-basis-and-fx-carry:3}
Swapping into dollars costs 40 basis points a year more than parity; borrowers of dollars through swaps (non-US banks, investors hedging dollar
assets) pay it to whoever can lend dollars through the swap and has the balance sheet.
\end{solution}

\begin{solution}{iq:s2:cross-currency-basis-and-fx-carry:4}
A 2008-style flight to safety (high-yielders down 15--20\%, funding currencies up), a rate shock in one funding currency, and a liquidity event that
widens forward spreads.
\end{solution}

\begin{solution}{iq:s2:cross-currency-basis-and-fx-carry:5}
Spot and forward points by tenor, deposit rates, holiday calendars by currency, value dates; month-end brings roll dates, holidays that shift value
dates and illiquid quotes.
\end{solution}

\begin{solution}{iq:s2:cross-currency-basis-and-fx-carry:6}
Borrowing one dollar at $r_\$$ or converting at spot $S$ (foreign per dollar), lending at $r_f$ and converting back at $F$ must give the same:
$F = S(1 + r_f\tau)/(1 + r_\$\tau)$. With a basis, $F = S(1 + (r_f + b)\tau)/(1 + r_\$\tau)$: a negative $b$ makes the foreign leg worth less, so dollars
obtained through the swap cost more.
\end{solution}
